\documentclass[12pt,a4paper]{article}
\usepackage[margin=2.3cm]{geometry}
\usepackage{xcolor}
\usepackage[T1]{fontenc}
\usepackage[utf8]{inputenc}
\usepackage[brazil]{babel}
\usepackage{lmodern}
\usepackage{amsmath,amssymb,booktabs,array}
\usepackage{tikz}
\usetikzlibrary{automata,positioning,arrows.meta}
\usepackage{listings}
\definecolor{azul}{RGB}{20,49,78}
\definecolor{verde}{RGB}{0,121,111}
\definecolor{cinza}{RGB}{244,247,250}
\lstset{language=Python,basicstyle=\ttfamily\scriptsize,
 keywordstyle=\color{azul}\bfseries,commentstyle=\color{verde},
 stringstyle=\color{verde},breaklines=true,showstringspaces=false,
 frame=single,rulecolor=\color{azul!20},backgroundcolor=\color{cinza},
 columns=fullflexible,keepspaces=true,tabsize=4}
\newcommand{\q}[1]{q_{\mathrm{#1}}}
\newcommand{\MiniL}{\textsc{MiniL}}
\newcommand{\grupo}{Gabriel Holanda, Gabriel Torres, Victor Milito,\\Paulo Júnior, Vinicius Kaua e Pablo Café}
\tikzset{every state/.style={draw=azul,thick,minimum size=8mm,font=\small},
 every edge/.style={draw=azul,>=Stealth},every edge quotes/.style={font=\scriptsize}}

\usepackage{microtype}
\usepackage[colorlinks=true,linkcolor=azul,urlcolor=verde]{hyperref}
\usepackage{fancyhdr}
\pagestyle{fancy}\fancyhf{}
\fancyhead[L]{\small\MiniL{} | Apresentação ao cliente}
\fancyfoot[C]{\thepage}
\setlength{\headheight}{15pt}
\setlength{\parindent}{0pt}\setlength{\parskip}{7pt}
\setlength{\emergencystretch}{2em}
\title{\color{azul}\textbf{MiniL: entendendo o texto antes de executar}\\[3mm]
\large Relatório explicativo para o cliente}
\author{\normalsize\grupo}
\date{Versão 1.1.0 --- 13 de novembro de 2026}
\begin{document}\maketitle
\section*{A proposta em uma página}
\textbf{O que entregamos.} Um programa que lê um arquivo de texto e identifica
suas unidades básicas: nomes, números, palavras reservadas, sinais de operação
e separadores. Cada unidade vem acompanhada de sua posição no arquivo. Quando
encontra um caractere que não pertence às regras combinadas, informa onde ele está.

\textbf{Para que serve.} Imagine uma equipe desenvolvendo um editor para uma
linguagem simples. Antes de conferir se os comandos estão bem organizados,
ela precisa reconhecer o que foi escrito. A MiniL entrega essa primeira
separação e permite que outro programa receba os resultados em formato JSON.
Esse cenário é ilustrativo: não alegamos uma implantação comercial ou um
ganho de produtividade medido.

\textbf{Como funciona para quem usa.} A pessoa entrega o arquivo; o programa
lê o conteúdo, separa as unidades e apresenta a lista. Se houver um erro,
a execução termina com uma mensagem de localização, sem apresentar uma lista
parcial como se fosse um resultado completo.

\begin{center}
\begin{tikzpicture}[>=Stealth,node distance=9mm,
 box/.style={draw=azul,rounded corners,fill=cinza,text width=35mm,
 align=center,minimum height=15mm,font=\small}]
\node[box] (a) {Arquivo de texto\\em UTF-8};
\node[box,right=of a] (b) {MiniL\\aplica as regras};
\node[box,right=of b] (c) {Lista com posições\\ou mensagem de erro};
\draw[->,thick,azul] (a)--(b);
\draw[->,thick,azul] (b)--(c);
\end{tikzpicture}
\end{center}

\textbf{Limite da entrega.} A MiniL não executa o arquivo nem verifica se o
programa está correto como um todo. Ela realiza a leitura inicial. A validação
da estrutura dos comandos e do significado das operações seria uma etapa futura.

\section*{Materiais entregues}
Este relatório explica o uso. O relatório técnico apresenta a especificação
matemática, os diagramas e a minimização. A pasta \texttt{slides} contém a
apresentação; \texttt{codigo}, \texttt{testes} e \texttt{exemplos} permitem
executar e conferir a solução. Todos os documentos e os slides possuem fontes LaTeX.

\newpage
\section{Um exemplo completo}
Para a entrada
\begin{lstlisting}[language={}]
let total = 12.50;
\end{lstlisting}
o cliente recebe:
\begin{center}
\begin{tabular}{lll}
\toprule Texto encontrado & Significado na MiniL & Linha:coluna\\\midrule
\texttt{let} & Palavra reservada (KW\_LET) & 1:1\\
\texttt{total} & Nome escolhido pelo usuário (ID) & 1:5\\
\texttt{=} & Operador (OP) & 1:11\\
\texttt{12.50} & Número decimal (FLOAT) & 1:13\\
\texttt{;} & Separador (SEP) & 1:18\\\bottomrule
\end{tabular}
\end{center}
Uma linha e uma coluna começam em 1. A posição permite localizar a unidade
no arquivo original. Espaços entre palavras são lidos, mas não entram nessa lista.

\section{Regras acordadas}
\begin{itemize}
\item Palavras reservadas: \texttt{let}, \texttt{if}, \texttt{else} e \texttt{print}.
\item Nomes: começam com letras de A a Z, maiúsculas ou minúsculas, ou com
sublinhado. Depois disso, podem conter dígitos. Acentos não são permitidos nos nomes.
\item Números: inteiros como \texttt{12}, e decimais como \texttt{12.50}.
O decimal exige dígitos antes e depois do ponto; vírgula não é aceita.
\item Operadores: \texttt{+}, \texttt{-}, \texttt{*}, \texttt{/} e \texttt{=}.
Separadores: \texttt{(}, \texttt{)} e \texttt{;}.
\item Maiúsculas e minúsculas são diferentes: \texttt{if} é reservada,
mas \texttt{If} é um nome comum.
\end{itemize}
O programa consome o maior trecho permitido antes de classificar. Por isso,
\texttt{ifx} é um único nome; \texttt{12abc} é separado em número \texttt{12}
e nome \texttt{abc}. O sinal de \texttt{-3} é um operador separado do número.
\texttt{==} representa dois operadores \texttt{=}; igualdade composta não foi definida.

\section{Como os erros são apresentados}
\begin{lstlisting}[language={}]
Entrada: let total = 12.;
Erro: Caractere inesperado '.' em 1:15
\end{lstlisting}
O ponto está sem a parte decimal. A pessoa deve corrigir o arquivo e executar
novamente. O programa para no primeiro erro; não tenta corrigir o conteúdo
automaticamente. Um arquivo inexistente ou sem codificação UTF-8 também recebe
mensagem de erro, sem uma listagem técnica de falha interna.

\newpage
\section{Uso e integração}
É necessário Python 3.12 ou superior. Não há bibliotecas externas a instalar.
Na pasta principal do projeto, os comandos básicos são:
\begin{lstlisting}[language={}]
python codigo/scanner.py exemplos/exemplo.min
python codigo/scanner.py exemplos/exemplo.min --json
python codigo/scanner.py exemplos/invalido.min
\end{lstlisting}
O primeiro mostra uma lista no terminal. O segundo apresenta os mesmos dados
em JSON, formato que um editor ou outra aplicação pode ler. O terceiro demonstra
um erro esperado. Uma execução bem-sucedida termina com código 0; erros terminam
com código 1. Erros de uso da linha de comando, como uma opção desconhecida,
terminam com código 2 e mostram a ajuda do Python.

\section{Como conferir a entrega}
\begin{center}\small
\begin{tabular}{>{\raggedright\arraybackslash}p{.35\linewidth}>{\raggedright\arraybackslash}p{.55\linewidth}}
\toprule Verificação do cliente & Resultado esperado\\\midrule
Ler \texttt{let total = 12.50;} & Cinco unidades, nas posições mostradas no exemplo\\
Ler \texttt{ifx} & Um nome comum; não separar a palavra \texttt{if}\\
Ler \texttt{let total = 12.;} & Mensagem de erro no ponto em 1:15, sem saída parcial\\
Usar a opção \texttt{--json} & Dados estruturados com categoria, texto, linha e coluna\\
Executar \texttt{python} com \path{scripts/validar.py} & 18 testes aprovados no ambiente registrado\\\bottomrule
\end{tabular}
\end{center}
Os testes cobrem exemplos positivos e negativos, posições, leitura de arquivos,
saída estruturada e propriedades do reconhecedor. O resultado reproduzível fica
em \path{evidencias/validacao.txt}; a versão do interpretador efetivamente usada
também está registrada. A revisão foi feita em novembro de 2026; isso não altera
as datas das técnicas teóricas que fundamentam o produto.

\section{O que fica fora do escopo}
Não há interface gráfica, execução de comandos, contas de usuário ou serviços
remotos. Textos entre aspas, notação científica e comentários não foram definidos.
Uma barra continua sendo um operador: duas barras não iniciam um comentário.
Parênteses são identificados individualmente; seu balanceamento não é verificado.
Um arquivo pode ser separado em unidades e ainda conter comandos inválidos.

Uma próxima versão poderia adicionar um editor visual e uma análise da
estrutura dos comandos. Essas funções são propostas de evolução, não recursos
da versão entregue.

\end{document}
